\documentclass[12pt]{article}

\usepackage{amsmath}
\usepackage[margin=1in]{geometry}
\usepackage{fancyhdr}
\usepackage{amsthm}
\usepackage{amssymb}
\usepackage{pgfplots}
\usepackage{tikz}
\usepackage{bbm}
\usepackage{bm}
\usepackage{dsfont}
\usetikzlibrary{shapes.geometric}
\usepackage{enumerate}
\usepackage[shortlabels]{enumitem}

\pgfplotsset{compat=1.18}

\newcommand{\pderiv}[2][]{\frac{\partial#1}{\partial#2}}

\pagestyle{fancy}
\fancyhead[l]{Hudson Benites}
\fancyhead[c]{Day 10}
\fancyfoot[c]{\thepage}
\renewcommand{\headrulewidth}{0.2pt}
\setlength{\headheight}{15pt}
\fancyhead[r]{\today}

\newcommand{\magnitude}[1]{\left\lvert #1 \right\rvert} %chktex 1

\begin{document}
\begin{enumerate}[label=\Alph*., leftmargin=*]
    \item \begin{enumerate}
        \item After playing around with $g$ and different choices of $w$ and $y$, I came to the conclusion that $g$ must be a scalar multiple of the the identity matrix $\mathds{1}$.
        \item \begin{proof}
            Let $g$ be a scalar multiple of the identity matrix. Let $w, y$ be vectors in $\mathbb{R}^2$. Then their dot product is $\frac{w^{T}\mathds{1}y}{\sqrt{w^{T}\mathds{1}w}\sqrt{y^{T}\mathds{1}y}}$. The angle between two vectors pushed forward by $dx_q$ is
            $\frac{w^{T}gy}{\sqrt{w^{T}gw}\sqrt{y^{T}gy}}$. When $g=c\mathds{1}$, this simplifies to $\frac{cw^{T}\mathds{1}y}{\sqrt{cw^{T}\mathds{1}w}\sqrt{cy^{T}\mathds{1}y}}=\frac{w^{T}\mathds{1}y}{\sqrt{w^{T}\mathds{1}w}\sqrt{y^{T}\mathds{1}y}}$. So, $dx_q$ is angle preserving.
        \end{proof}
    \end{enumerate}
    \item \begin{enumerate}
        \item The $g$ associated with $x(u,v)=(\frac{2u}{u^2+v^2+1}, \frac{2v}{u^2+v^2+1}, \frac{u^2+v^2-1}{u^2+v^2+1})$ is $g_x=\frac{4}{{(u^2+v^2+1)}^2}\mathds{1}$. No matter what $q$ is selected,
        $g_x$ is always a scalar multiple of the identity, and thus preserves angles.
        \item When $q=(0,0)$, $g_x$ is $4\mathds{1}$, and thus $\sqrt{\det g}=4$ and $dx_q$ does not preserve area at this point.
        When $q$ is a point on the unit circle, $u^2+v^2=1$, and $g_x=\mathds{1}$, whose determinant is 1. Therefore, on the unit circle, $dx_q$ preserves area.
        As $q$ gets arbitrarily far away from the origin, $g$ approaches the 0 matrix, and thus the determinant is 0, which does not preserve area.
    \end{enumerate}
    \item Let $U=(-\pi, \pi)\times (-1, 1)$. We can send $U$ to $T$ by $r: U\to T$ given by $r(u_1,u_2)=(\cos u_1, \sin u_1, u_2)$.
    To project this onto the unit sphere, we need to make sure that $\magnitude{r(u,v)}^2$ is 1, but we also must keep the height the same, thus fixing $r_3$. So, we compute $\magnitude{r(u,v)}^2=\cos^2 u +\sin^2 u + v^2$. To make this 1,
    we can have $x(u,v)=(\sqrt{1-v^2}\cos u, \sqrt{1-v^2}\sin u, v)$. This is a local parameterization of $\mathbb{S}^2$. We now show that it is area preserving. We compute that $\pderiv[x]{u}=(-\sqrt{1-v^2}\sin u, \sqrt{1-v^2}\cos u, 0)$. We also compute that $\pderiv[x]{v}=(\frac{-v\cos u}{\sqrt{1-v^2}}, \frac{-v\sin u}{\sqrt{1-v^2}}, 1)$.
    We compute that $g=\begin{bmatrix}
        1-v^2&0\\
        0&\frac{1}{1-v^2}
    \end{bmatrix}$. We can see that $\det g = 1$, and thus $\sqrt{\det g}= 1$, and $x$ is area preserving.
    \end{enumerate}
\end{document}